\section{Выводы}

Выполнив первую лабораторную работу по курсу \enquote{Дискретный анализ}, я изучил принципы работы линейных алгоритмов сортировки и реализовал стабильный алгоритм поразрядной сортировки (Radix Sort LSD) в связке с сортировкой подсчетом.

В процессе выполнения лабораторной работы я:
\begin{itemize}
    \item Углубил знания языка C++, реализовав собственный контейнер динамического массива (\texttt{Vector}), аналогичный \texttt{std::vector}.
    \item Научился применять семантику перемещения (\texttt{std::move}) для оптимизации копирования тяжелых объектов (строк), что критически важно для производительности сортировки.
    \item Разобрался с нюансами быстрого парсинга данных и оптимизацией стандартных потоков ввода-вывода (\texttt{std::cin.tie}, \texttt{std::ios\_base::sync\_with\_stdio}).
    \item Написал Python-скрипты для генерации тестов и автоматической проверки (чекер) правильности сортировки и целостности данных.
    \item Освоил систему автоматизации сборки \texttt{make}.
\end{itemize}

Результаты бенчмарков наглядно показали преимущество специализированных линейных алгоритмов над универсальными сортировками сравнением при больших объемах данных.
\pagebreak